% ============================================================================
% A complete proof of \Cref{lemma:ExtensionOfTamelyRamifiedSheaf}.
%
% This file is intended to REPLACE the placeholder proof in
%   sections/GeometricCFT/ProofOfReducedTheorem.tex (the lemma at lines 42-53,
%   whose proof body is currently a red note).
% The lemma is restated below with the SAME label
%   \label{lemma:ExtensionOfTamelyRamifiedSheaf}
% so all existing \Cref's continue to work. When integrating, delete the old
% lemma+proof from ProofOfReducedTheorem.tex and \input this file in its place
% (otherwise the label will be multiply defined).
% ============================================================================

We keep the notation of \Cref{section:BlowupOfSymmetricPowerOfCurves}:
$C$ is a projective smooth geometrically connected curve over the perfect field $k$,
$\mdls = \sum n_i P_i > 0$ is a modulus with complement $U = C \setminus \mdls$,
and $d$ is an integer satisfying \Cref{equation:degree-condition}.
We write $Z_0 = Z_0(\mdls, d) \subset C^{(d)}$ for the center of the blowup
$\pi : X_{\mdls, d} \to C^{(d)}$, $E_0$ for its exceptional divisor with generic point $\eta_0$,
$\Cmod{d} \subseteq X_{\mdls,d}$ for Takeuchi's compactification
(\Cref{theorem:takeuchi-compactification-theorem}), and
$H = \Cmod{d} \setminus U^{(d)} = E_0 \times_{C^{(d)}} \Cmod{d}$ for the boundary divisor.
Throughout, $p$ denotes the characteristic of $k$ (the case $p=0$ only simplifies the arguments below).

\begin{lemma}\label{lemma:ExtensionOfTamelyRamifiedSheaf}
    If $\cF^{[d]}$ is a locally constant sheaf on $U^{(d)}$ which is tamely ramified along the boundary divisor $H = \tilde{C}^{(d)}_\mdls \setminus U^{(d)}$,
    then $\cF^{[d]}$ extends to a locally constant sheaf $\tilde{\cF}^{[d]}$ on $\tilde{C}^{(d)}_\mdls$.
\end{lemma}

The proof occupies the rest of this subsection and proceeds in two independent steps.
First (\Cref{prop:extension_via_purity}), by the Zariski--Nagata purity of the branch locus,
extending $\cF^{[d]}$ across $H$ is possible as soon as the corresponding torsor is
\emph{unramified} at the generic point $\eta_0$ of $H$; a priori tameness is weaker than this.
Second (\Cref{prop:tame_at_boundary_implies_unramified}), we show that at $\eta_0$
tameness in fact forces unramifiedness. This is where the global geometry enters:
the wild ($p$-primary) part of the torsor is killed at $\eta_0$ by tameness alone,
while the tame (prime-to-$p$) part is killed by an explicit Kummer-theoretic computation of valuations,
using the projective bundle structure of $C^{(d)}$ over $\PicC[d]$.
Note that the second step is exactly where ``tame'' is upgraded to ``unramified'':
for a \emph{sub}modulus $\ndls \subsetneq \mdls$ the analogous statement at $\eta_\ndls$ is false in general
(the computations of \Cref{section:ramification_after_blowup_is_tame} produce genuinely, tamely, ramified torsors there),
and the argument below makes visible why the full modulus and the completed linear systems are needed.

We first record the relevant geometry of the pair $(\Cmod{d}, H)$.

\begin{lemma}\label{lemma:boundary_of_compactification}
    With notation as above:
    \begin{enumerate}
        \item $\Cmod{d}$ is an integral scheme, smooth over $k$.
        \item $H$ is a nonempty open subscheme of $E_0$. In particular $H$ is irreducible of
              codimension $1$ in $\Cmod{d}$, its generic point is $\eta_0$, and $\eta_0 \in \Cmod{d}$.
        \item $\eta_0$ is the unique point of $\Cmod{d} \setminus U^{(d)}$ whose local ring has dimension $1$,
              and $\Oo_{\Cmod{d}, \eta_0} = \Oo_{X_{\mdls,d}, \eta_0}$ is a discrete valuation ring
              with fraction field $K := k(C^{(d)})$.
    \end{enumerate}
\end{lemma}
\begin{proof}
    (1) $\Cmod{d}$ is an open subscheme of $X_{\mdls,d}$, which is integral as the blowup of the
    integral scheme $C^{(d)}$; hence $\Cmod{d}$ is integral.
    By \Cref{theorem:takeuchi-compactification-theorem}, $\Cmod{d}$ is a projective space bundle over
    $\PicCm[d]$. The latter is smooth over $k$: $\PicCm[0]$ is a smooth group scheme and $\PicCm[d]$
    is a torsor under it (see the properties of the generalized Picard scheme recalled above).
    A projective space bundle over a smooth $k$-scheme is smooth over $k$.

    (2) Since $\Cmod{d} \hookrightarrow X_{\mdls,d}$ is an open immersion, its base change
    $H = E_0 \times_{C^{(d)}} \Cmod{d} = E_0 \times_{X_{\mdls,d}} \Cmod{d} \hookrightarrow E_0$
    is an open immersion as well. To see $H \neq \emptyset$, suppose otherwise; then
    $U^{(d)} = \Cmod{d}$ and $\Phi_d = \tilde{\Phi}_d$ would be a projective space bundle. But the
    fibers of $\Phi_d$ are affine spaces of dimension $d - \deg \mdls - g + 1 \geq 1$
    (\Cref{section:AbelJacobi}; positivity follows from \Cref{equation:degree-condition}),
    and a positive dimensional affine space over a field is not proper --- a contradiction.
    A nonempty open subscheme of the irreducible scheme $E_0$ is irreducible, dense, and contains the
    generic point $\eta_0$ of $E_0$.

    (3) By \Cref{theorem:takeuchi-compactification-theorem}, $\Cmod{d} \setminus U^{(d)} = H$,
    which by (2) is irreducible of codimension $1$ with generic point $\eta_0$; every other point of $H$
    is a proper specialization of $\eta_0$ and hence has local ring of dimension $\geq 2$.
    By (1), $\Oo_{\Cmod{d}, \eta_0}$ is a regular local ring of dimension $1$, i.e.\ a discrete
    valuation ring, and it equals $\Oo_{X_{\mdls,d}, \eta_0}$ because $\Cmod{d}$ is open in $X_{\mdls,d}$.
    Its fraction field is the function field of $X_{\mdls,d}$, which equals $K = k(C^{(d)})$
    since $\pi$ is birational.
\end{proof}

We denote by $v_{\eta_0}$ the discrete valuation on $K$ associated to $\Oo_{\Cmod{d}, \eta_0}$,
and by $\widehat{K}_{\eta_0}$ the fraction field of the completion.

\subsubsection*{Step 1: Purity, and extending unramified torsors}

\begin{prop}\label{prop:extension_via_purity}
    Let $G$ be a finite abelian group, let $\rho : \pi_1^{et}(U^{(d)}, \bar{u}) \to G$ be a continuous
    homomorphism and let $\cQ$ be the corresponding $G$-torsor on $U^{(d)}_{\text{\'et}}$
    (\Cref{cor:abelian_equivilance_fundamental_torsors}).
    If $\cQ$ is unramified over $\Cmod{d}$ at $(H, \eta_0)$ in the sense of
    \Cref{definition:tamely_ramified_in_codim_1}, then $\rho$ factors, necessarily uniquely, through the
    canonical surjection
    \[ \iota_* : \pi_1^{et}(U^{(d)}, \bar{u}) \twoheadrightarrow \pi_1^{et}(\Cmod{d}, \bar{u}), \]
    where $\iota : U^{(d)} \hookrightarrow \Cmod{d}$ is the open immersion.
    Consequently, $\cQ$ extends to a $G$-torsor on $(\Cmod{d})_{\text{\'et}}$, uniquely up to isomorphism.
\end{prop}
\begin{proof}
    \textbf{Surjectivity of $\iota_*$.}
    We first note that for every connected finite \'etale cover $T \to \Cmod{d}$, the restriction
    $T \times_{\Cmod{d}} U^{(d)}$ is connected: $T$ is \'etale over the regular scheme $\Cmod{d}$, hence
    regular, and being connected it is irreducible; $T \times_{\Cmod{d}} U^{(d)}$ is a nonempty open
    subscheme of $T$ (nonempty since $T \to \Cmod{d}$ is surjective and $U^{(d)}$ is dense), hence
    irreducible, hence connected.
    This implies surjectivity of $\iota_*$: if the image were a proper closed subgroup, it would be
    contained in a proper open subgroup $M \leq \pi_1^{et}(\Cmod{d}, \bar{u})$, and the connected cover
    corresponding to $M$ would pull back to a disconnected cover of $U^{(d)}$ (the image of
    $\pi_1^{et}(U^{(d)})$ does not act transitively on $\pi_1^{et}(\Cmod{d})/M$), a contradiction.

    \textbf{The connected cover and its normalization.}
    Let $G' = \rho\left(\pi_1^{et}(U^{(d)}, \bar{u})\right) \subseteq G$ and let $V \to U^{(d)}$ be the connected
    finite \'etale Galois cover with group $G'$ corresponding to the open normal subgroup
    $\ker \rho$ (\Cref{theorem:equivalence_fundamental_covers}). As a $U^{(d)}$-scheme, $\cQ$ is
    isomorphic to a disjoint union of $[G : G']$ copies of $V$; hence at the generic point $\eta_0$ the
    finite separable $\widehat{K}_{\eta_0}$-algebras of $\cQ$ and of $V$ have the same field factors, and
    $V$ is unramified over $\Cmod{d}$ at $(H, \eta_0)$.

    $V$ is \'etale over the smooth scheme $U^{(d)}$, hence smooth over $k$; being connected it is
    integral, and $K(V)/K$ is a finite separable field extension. Let
    $\nu : W \to \Cmod{d}$ be the normalization of $\Cmod{d}$ in $K(V)$, i.e.\ the relative spectrum of
    the integral closure of $\Oo_{\Cmod{d}}$ in $K(V)$. Since $\Cmod{d}$ is of finite type over a field it
    is Nagata (\stackstag{0335}), so $\nu$ is a finite morphism, and $W$ is integral and normal.
    Over the open $U^{(d)}$, the scheme $V$ is finite over $U^{(d)}$, normal, with function field $K(V)$;
    by uniqueness of normalization, $\nu^{-1}(U^{(d)}) \cong V$ as $U^{(d)}$-schemes.

    \textbf{$\nu$ is \'etale, by purity of the branch locus.}
    We apply \stackstag{0BMB} at each point $w \in W$ with image $y = \nu(w)$. The hypotheses hold:
    $\Oo_{W,w}$ is normal; $\Oo_{\Cmod{d},y}$ is regular (\Cref{lemma:boundary_of_compactification});
    $\nu$ is finite, hence quasi-finite at $w$; and
    $\dim(\Oo_{W,w}) = \dim(\Oo_{\Cmod{d},y})$, since for a finite dominant morphism of integral schemes
    of finite type over a field one has
    $\dim(\Oo_{W,w}) = \dim W - \dim \overline{\{w\}}$, $\dim W = \dim \Cmod{d}$ and
    $\dim \overline{\{w\}} = \dim \overline{\{y\}}$.
    It remains to check the codimension-one condition: let $w' \rightsquigarrow w$ be a specialization
    with $\dim(\Oo_{W, w'}) = 1$; then $y' = \nu(w')$ satisfies $\dim(\Oo_{\Cmod{d}, y'}) = 1$
    by the same dimension count. If $y' \in U^{(d)}$, then $\nu$ is \'etale at $w'$ since
    $\nu^{-1}(U^{(d)}) = V$ is \'etale over $U^{(d)}$. Otherwise
    $y' \in \Cmod{d} \setminus U^{(d)}$, and by \Cref{lemma:boundary_of_compactification}(3) necessarily
    $y' = \eta_0$; then $\Oo_{W, w'}$ is a localization of the integral closure of the discrete valuation
    ring $\Oo_{\Cmod{d},\eta_0}$ in $K(V)$ at one of its maximal ideals, and the hypothesis that $V$ is
    unramified at $(H, \eta_0)$ says precisely (\Cref{definition:tamely_ramified_dvrs}) that the
    ramification index there is $1$ and the residue extension is separable, i.e.\ that $\nu$ is
    unramified at $w'$. If $\dim(\Oo_{W,w}) = 0$ then $w$ is the generic point of $W$ and $\nu$ is \'etale
    at $w$ because $K(V)/K$ is finite separable. In all cases \stackstag{0BMB} applies, and we conclude
    that $\nu : W \to \Cmod{d}$ is finite \'etale.

    \textbf{Factorization.}
    The geometric point $\bar{u}$ of $U^{(d)}$ is also a geometric point of $\Cmod{d}$, and
    $W_{\bar{u}} = V_{\bar{u}}$ as finite sets. The action of $\pi_1^{et}(U^{(d)}, \bar u)$ on
    $V_{\bar u}$ is the restriction along $\iota_*$ of the action of $\pi_1^{et}(\Cmod{d}, \bar u)$ on
    $W_{\bar u}$, because $V = W \times_{\Cmod{d}} U^{(d)}$. Since $V$ is Galois with group $G'$, the
    kernel of the first action is exactly $\ker \rho$; hence
    $\ker(\iota_*) \subseteq \ker \rho$. As $\iota_*$ is surjective, $\rho$ factors as
    $\rho = \tilde{\rho} \circ \iota_*$ for a unique continuous
    $\tilde{\rho} : \pi_1^{et}(\Cmod{d}, \bar u) \to G' \subseteq G$.
    Let $\tilde{\cQ}$ be the $G$-torsor on $(\Cmod{d})_{\text{\'et}}$ corresponding to $\tilde{\rho}$
    (\Cref{cor:abelian_equivilance_fundamental_torsors}); then
    $\iota^* \tilde{\cQ} \cong \cQ$ since the corresponding characters of $\pi_1^{et}(U^{(d)}, \bar u)$
    agree. Finally, if $\tilde{\cQ}_1, \tilde{\cQ}_2$ are two extensions of $\cQ$, their characters
    $\tilde{\rho}_1, \tilde{\rho}_2$ satisfy $\tilde{\rho}_1 \circ \iota_* = \tilde{\rho}_2 \circ \iota_*$,
    hence are equal by surjectivity of $\iota_*$, and $\tilde{\cQ}_1 \cong \tilde{\cQ}_2$.
\end{proof}

\subsubsection*{Step 2: At the boundary, tame implies unramified}

\begin{prop}\label{prop:tame_at_boundary_implies_unramified}
    Let $G$ be a finite abelian group and let $\cQ$ be a $G$-torsor on $U^{(d)}_{\text{\'et}}$.
    If $\cQ$ is tamely ramified over $\Cmod{d}$ at $(H, \eta_0)$, then $\cQ$ is unramified at
    $(H, \eta_0)$.
\end{prop}

Note that no bound on the ramification of $\cQ$ along $\mdls$ is assumed;
being a torsor on all of $U^{(d)}$ is enough.

\begin{proof}
    As in \Cref{section:ramification_after_blowup_is_tame} (see the discussion following
    \Cref{definition:local_ramification_boundness}), the restriction of $\cQ$ to
    $\Spec(\widehat{K}_{\eta_0})$ decomposes as a disjoint union of copies of $\Spec(F)$ for a finite
    Galois extension $F / \widehat{K}_{\eta_0}$ with Galois group a subgroup of $G$, and corresponds to a
    continuous homomorphism
    \[ \rho : G_{\widehat{K}_{\eta_0}} \to G \]
    with image $\Gal(F/\widehat{K}_{\eta_0})$. Let $I \subseteq G_{\widehat{K}_{\eta_0}}$ be the inertia
    subgroup. Then $\cQ$ is unramified at $(H,\eta_0)$ if and only if $\rho(I) = \{1\}$, and
    $\rho(I)$ equals the inertia subgroup $G_0$ of $\Gal(F/\widehat{K}_{\eta_0})$.

    \textbf{Reduction to $k$ algebraically closed.}
    Let $k'/k$ be a finite separable extension. The projection
    $\mathrm{pr} : (\Cmod{d})_{k'} \to \Cmod{d}$ is (finite) \'etale, and it identifies
    $(\Cmod{d})_{k'} \setminus (U^{(d)})_{k'}$ with $\mathrm{pr}^{-1}(H)$. Since $C^{(n)}$ is
    geometrically integral for every $n$ (see the Preliminaries), so are $Z_0$, $C^{(d)}$ and hence
    $E_0$ (the exceptional divisor of a blowup of a geometrically integral smooth scheme along a
    geometrically integral smooth center); therefore $\mathrm{pr}^{-1}(H)$ is irreducible with a unique
    generic point $\eta'$ lying over $\eta_0$, and $\mathrm{pr}$ is \'etale at $\eta'$. By
    \Cref{lemma:descend_ramification_along_etale} (and its proof: the completed local extensions at
    $\eta'$ and at $\eta_0$ differ by the unramified base change
    $\widehat{K}_{\eta'} / \widehat{K}_{\eta_0}$), the torsor $\mathrm{pr}^{-1}\cQ$ is tamely ramified,
    resp.\ unramified, at $\eta'$ if and only if $\cQ$ is tamely ramified, resp.\ unramified, at
    $\eta_0$. Passing to the colimit over all finite separable $k'/k$ (recall $k$ is perfect, so
    $\bar{k} = k^{sep}$), we may and do assume for the rest of the proof that $k$ is algebraically
    closed. In particular the points of $\mdls$ are $k$-rational; write
    $\operatorname{supp}(\mdls) = \{P_1, \dots, P_s\}$.

    \textbf{Splitting off the wild part.}
    By the structure theorem for finite abelian groups, write $G = G_p \times G'$ with $G_p$ a $p$-group
    and $|G'|$ prime to $p$ (if $p = 0$, take $G_p = \{1\}$). Then
    $\rho(I) = \{1\}$ if and only if both compositions
    $\rho_p := \mathrm{pr}_{G_p} \circ \rho$ and $\rho' := \mathrm{pr}_{G'} \circ \rho$ kill $I$.

    Since $\cQ$ is tamely ramified at $(H, \eta_0)$, the extension $F/\widehat{K}_{\eta_0}$ is tamely
    ramified with respect to $\Oo_{\Cmod{d}, \eta_0}$ (\Cref{definition:tamely_ramified_dvrs}): its
    residue extension is separable and its ramification index $e$ is prime to $p$. For a Galois
    extension of complete discrete valuation rings with separable residue extension, the inertia group
    $G_0$ has order equal to the ramification index (\cite{serreLF}, Ch.\ IV); hence
    $\rho(I) = G_0$ is a group of order prime to $p$. Its projection to the $p$-group $G_p$ is then
    trivial, that is:
    \[ \rho_p(I) = \{1\}, \qquad \rho(I) \subseteq \{1\} \times G'. \]
    It therefore remains to prove $\rho'(I) = \{1\}$.

    \textbf{The prime-to-$p$ part, via Kummer theory.}
    Let $N = |G'|$, which is invertible in $k = \bar{k}$, so $\mu_N \subset k$. Since $G'$ is abelian,
    $\rho'(I) = \{1\}$ holds if and only if $\chi(\rho'(I)) = \{1\}$ for every character
    $\chi : G' \to \mu_N(k)$. Fix such a $\chi$ and let $e \mid N$ be the order of the character
    \[ \chi \circ \rho'_{glob} : \pi_1^{et}(U^{(d)}, \bar{u}) \longrightarrow \mu_e(k), \]
    where $\rho'_{glob}$ denotes the global character of $\cQ / G_p$ (\Cref{prop:quotient_torsors})
    whose local restriction at $\eta_0$ is $\rho'$. It suffices to show that the $\mu_e$-torsor
    corresponding to $\chi \circ \rho'_{glob}$ is unramified at $(H, \eta_0)$; indeed, the inertia $I$
    then acts trivially on it, i.e.\ $\chi(\rho'(I)) = \{1\}$.

    Restrict $\chi \circ \rho'_{glob}$ to the absolute Galois group of the function field
    $K = k(U^{(d)}) = k(C^{(d)})$. Since $\mu_e \subset K$, Kummer theory identifies
    $H^1(K, \mu_e) \cong K^\times / (K^\times)^e$; let $F \in K^\times$ be a representative of the
    resulting class, so that the corresponding cyclic cover of $\Spec(K)$ is
    $\Spec\left(K[y]/(y^e - F)\right)$, and $F$ is well defined up to $e$-th powers. Accordingly,
    $\operatorname{div}_{C^{(d)}}(F)$ is well defined modulo $e \cdot \operatorname{Div}(C^{(d)})$.
    We will show:
    \begin{equation}\label{equation:kummer_valuation_at_boundary}
        v_{\eta_0}(F) \equiv 0 \pmod{e}.
    \end{equation}
    Let us first explain why \eqref{equation:kummer_valuation_at_boundary} finishes the proof.
    Let $\varpi$ be a uniformizer of $B_0 := \Oo_{\Cmod{d}, \eta_0}$ and write
    $v_{\eta_0}(F) = e a$; replacing $F$ by $F \varpi^{-ea}$ (an allowed change modulo nothing --- it
    changes neither the class in $H^1(\widehat{K}_{\eta_0}, \mu_e)$ nor the extension, since it changes
    $F$ by the $e$-th power $\left(\varpi^{-a}\right)^e$), we may assume $F \in \widehat{B}_0^\times$, where
    $\widehat{B}_0$ is the completion. Then
    \[ B := \widehat{B}_0[y]/(y^e - F) \]
    is finite \'etale over $\widehat{B}_0$: it is finite free, and the derivative $e y^{e-1}$ is a unit
    in $B$ because $e$ is invertible in $k$ and $y^e = F$ is a unit. Consequently
    $\widehat{K}_{\eta_0}[y]/(y^e - F) = B \otimes_{\widehat{B}_0} \widehat{K}_{\eta_0}$ is a product of
    finite separable field extensions of $\widehat{K}_{\eta_0}$, each unramified with respect to
    $\widehat{B}_0$ (\Cref{definition:tamely_ramified_dvrs}; compare
    \Cref{theorem:kummer_ramification}(1),(2), whose proof for $e$ prime to $p$ does not use perfectness
    of the residue field). Thus the $\mu_e$-torsor of $\chi \circ \rho'_{glob}$ is unramified at
    $(H, \eta_0)$, as desired.

    The rest of the proof establishes \eqref{equation:kummer_valuation_at_boundary}, in three claims.
    For each $1 \leq i \leq s$ let
    \[ \Delta_{P_i} := Z_0(P_i, d) \subset C^{(d)} \]
    be the prime divisor obtained as the (closed) image of the morphism
    $C^{(d-1)} \to C^{(d)}$, $D \mapsto D + P_i$; thus
    $C^{(d)} \setminus U^{(d)} = \bigcup_{i=1}^{s} \Delta_{P_i}$, because a divisor of degree $d$ fails
    to be prime to $\mdls$ exactly when it contains some $P_i$. Set
    \[ c_i := v_{\Delta_{P_i}}(F) \in \Z. \]

    \begin{claim}\label{claim:divisor_congruence}
        $\operatorname{div}_{C^{(d)}}(F) \equiv \sum_{i=1}^{s} c_i \, \Delta_{P_i}
        \pmod{e \cdot \operatorname{Div}(C^{(d)})}$.
    \end{claim}
    \begin{proof}
        Let $W \subset C^{(d)}$ be a prime divisor different from all the $\Delta_{P_i}$, with generic
        point $\xi_W$. Then $\xi_W \in U^{(d)}$, and the $\mu_e$-torsor of $\chi \circ \rho'_{glob}$,
        being a torsor on all of $U^{(d)}$, is finite \'etale over a neighborhood of $\xi_W$, hence
        unramified at $(W, \xi_W)$ (\Cref{definition:tamely_ramified_in_codim_1}, with ambient scheme
        $C^{(d)}$). On the other hand, if $e \nmid v_W(F)$, then in any extension of the valuation $v_W$
        to $K(y)/(y^e - F)$ one has $v(y) = v_W(F)/e \notin v_W(K^\times)$, so the ramification index at
        $(W, \xi_W)$ is $e / \gcd(e, v_W(F)) > 1$ --- contradicting unramifiedness. Hence
        $e \mid v_W(F)$ for every such $W$.
    \end{proof}

    Let $\zeta$ denote the generic point of $Z_0$ and let $A := \Oo_{C^{(d)}, \zeta}$, a regular local
    ring of dimension $\deg \mdls$ (the codimension of $Z_0 \cong C^{(d - \deg\mdls)}$ in $C^{(d)}$)
    with maximal ideal $\mdls_\zeta = I_{Z_0, \zeta}$ and residue field $\kappa(\zeta)$.
    Since $E_0$ dominates $Z_0$, the center of $v_{\eta_0}$ on $C^{(d)}$ is $\zeta$, i.e.\
    $A \subseteq \Oo_{X_{\mdls,d},\eta_0}$ with $\mdls_\zeta = A \cap \mathfrak{m}_{\eta_0}$;
    in particular $v_{\eta_0} \geq 0$ on $A$ and $v_{\eta_0} = 0$ on $A^\times$.
    For $0 \neq g \in A$ let
    $\operatorname{ord}_\zeta(g) := \max\{ n : g \in \mdls_\zeta^n \}$
    be the order of vanishing of $g$ along $Z_0$.

    \begin{claim}\label{claim:exceptional_valuation_is_order}
        For every $0 \neq g \in A$ we have $v_{\eta_0}(g) = \operatorname{ord}_\zeta(g)$.
    \end{claim}
    \begin{proof}
        This is the general form of the local computation carried out explicitly (in coordinates) in
        \Cref{section:ramification_after_blowup_is_tame}; we recall the argument.
        Since $\Spec(A) \to C^{(d)}$ is flat and $I_{Z_0} \cdot A = \mdls_\zeta$, blowups commute with
        this base change (\stackstag{0805}):
        \[ X_{\mdls, d} \times_{C^{(d)}} \Spec(A) \cong \Bl_{\mdls_\zeta}(\Spec A), \]
        the exceptional divisor pulls back to the exceptional divisor
        $E_A = \mathbb{P}^{\deg\mdls - 1}_{\kappa(\zeta)}$, whose generic point maps to $\eta_0$ with
        the same local ring. So we may compute $v_{\eta_0}$ on the blowup of the regular local ring $A$
        at its maximal ideal. Choose a regular system of parameters $u_1, \dots, u_n$ of $A$
        ($n = \deg \mdls$) and consider the chart $A' = A[\mdls_\zeta / u_n]$. One has
        $\mdls_\zeta A' = (u_n)$ and, since $A$ is regular,
        $A' / (u_n) \cong \kappa(\zeta)[\bar{v}_1, \dots, \bar{v}_{n-1}]$ is a polynomial ring in the
        classes of $v_j := u_j / u_n$; in particular $(u_n) \subset A'$ is prime, the localization
        $A'_{(u_n)} = \Oo_{X_{\mdls,d}, \eta_0}$ is a discrete valuation ring with uniformizer $u_n$,
        and $E_A$ meets this chart in its generic point.
        Now let $r = \operatorname{ord}_\zeta(g)$ and write $g = G_r(u_1, \dots, u_n) + g'$ where $G_r$
        is a homogeneous form of degree $r$ with unit or zero coefficients, nonzero as an element of
        $\operatorname{gr}^r_{\mdls_\zeta}(A) \cong \kappa(\zeta)[\bar{u}_1, \dots, \bar{u}_n]_r$, and
        $g' \in \mdls_\zeta^{r+1}$. In the chart,
        \[ g = u_n^{\,r}\left( G_r(v_1, \dots, v_{n-1}, 1) + u_n h \right), \qquad h \in A', \]
        and the residue of $G_r(v_1, \dots, v_{n-1}, 1)$ in $A'/(u_n)$ is the dehomogenization of the
        nonzero form $G_r$, hence nonzero; therefore $G_r(v, 1) + u_n h$ is a unit in $A'_{(u_n)}$ and
        $v_{\eta_0}(g) = r$.
    \end{proof}

    \begin{claim}\label{claim:boundary_divisors_have_multiplicity_one}
        For each $i$, every local equation $w_i \in A$ of $\Delta_{P_i}$ at $\zeta$ satisfies
        $\operatorname{ord}_\zeta(w_i) = 1$; consequently $v_{\eta_0}(w_i) = 1$.
    \end{claim}
    \begin{proof}
        First note $\zeta \in \Delta_{P_i}$: every divisor $D \geq \mdls$ contains $P_i$, i.e.\
        $Z_0 \subseteq \Delta_{P_i}$. Since $A$ is regular local, hence a unique factorization domain,
        the height one prime of $\Delta_{P_i}$ at $\zeta$ is principal; fix a generator $w_i$.

        By \Cref{cor:etale_at_generic_point}, the addition morphism
        \[ \sigma : C^{(\deg\mdls)} \times_k C^{(d - \deg\mdls)} \longrightarrow C^{(d)} \]
        is \'etale at the point $\zeta' := (z_\mdls, \xi')$, where $z_\mdls \in C^{(\deg\mdls)}$ is the
        closed point corresponding to $\mdls$ and $\xi'$ is the generic point of $C^{(d-\deg\mdls)}$;
        moreover $\sigma(\zeta') = \zeta$, since $\sigma$ restricted to
        $\{z_\mdls\} \times C^{(d-\deg\mdls)}$ is the isomorphism onto $Z_0$ defining it.
        Because $\sigma$ is \'etale at $\zeta'$, the local homomorphism
        $A \to A'' := \Oo_{C^{(\deg\mdls)} \times C^{(d-\deg\mdls)},\, \zeta'}$ is faithfully flat with
        $\mdls_\zeta A'' = \mathfrak{m}_{\zeta'}$, so $\operatorname{ord}_\zeta(g) =
        \operatorname{ord}_{\zeta'}(\sigma^* g)$ for all $g \in A$.

        We identify $\sigma^* w_i$ near $\zeta'$. Set-theoretically,
        \[ \sigma^{-1}(\Delta_{P_i}) =
        \left( \Delta'_{P_i} \times_k C^{(d - \deg\mdls)} \right) \cup
        \left( C^{(\deg\mdls)} \times_k \Delta''_{P_i} \right), \]
        where $\Delta'_{P_i} \subset C^{(\deg\mdls)}$ and $\Delta''_{P_i} \subset C^{(d-\deg\mdls)}$ are
        the analogous divisors of divisors containing $P_i$. The second component does not pass through
        $\zeta'$ (the generic divisor of degree $d - \deg\mdls$ does not contain $P_i$), so near
        $\zeta'$ the scheme-theoretic preimage $\sigma^{-1}(\Delta_{P_i}) = V(\sigma^* w_i)$ is
        supported on the prime divisor $\Delta'_{P_i} \times_k C^{(d-\deg\mdls)}$ (a product of
        integral schemes, integral by \Cref{theorem:int_geoint_implies_productint}). Moreover
        $\sigma^{-1}(\Delta_{P_i})$ is reduced, being the \'etale pullback (near $\zeta'$) of the reduced
        scheme $\Delta_{P_i}$. Hence $\sigma^* w_i$ is, up to a unit at $\zeta'$, a local equation
        $w'_i \otimes 1$ of $\Delta'_{P_i} \times_k C^{(d-\deg\mdls)}$, where $w'_i$ is a local equation
        of $\Delta'_{P_i}$ at $z_\mdls$. Writing $B := \Oo_{C^{(\deg\mdls)}, z_\mdls}$ and noting that
        $A''$ is a localization of $B \otimes_k \kappa(\xi')$ at the prime $\mathfrak{m}_B \otimes \kappa(\xi')$,
        the homomorphism $B \to A''$ is flat with $\mathfrak{m}_B A'' = \mathfrak{m}_{\zeta'}$, so
        \[ \operatorname{ord}_\zeta(w_i) = \operatorname{ord}_{\zeta'}(w'_i \otimes 1)
           = \operatorname{ord}_{z_\mdls}(w'_i), \]
        and we are reduced to showing that $w'_i \notin \mathfrak{m}_{z_\mdls}^2$.

        This is a computation at the closed point $z_\mdls = \sum_j n_j P_j$ of $C^{(\deg\mdls)}$,
        for which we use the local description of symmetric powers from the Preliminaries
        (proof of the smoothness of $C^{(n)}$): choosing pairwise disjoint opens $U_j \ni P_j$ and
        local coordinates $t_j$ on $U_j$ vanishing exactly at $P_j$, a Zariski-open neighborhood of
        $z_\mdls$ decomposes as $\prod_j U_j^{(n_j)}$, and the completed local ring is
        \[ \widehat{\Oo}_{C^{(\deg\mdls)}, z_\mdls} \cong
           k[[\, \varepsilon^{(j)}_l \; : \; 1 \leq j \leq s, \; 1 \leq l \leq n_j \,]], \]
        where $\varepsilon^{(j)}_l$ is the $l$-th elementary symmetric function of the coordinates
        $t_j(x_1), \dots, t_j(x_{n_j})$ of the points of the $j$-th block. Under the product
        decomposition, $\Delta'_{P_i}$ is the preimage of the divisor
        $\{D_i : D_i \ni P_i\} \subset U_i^{(n_i)}$, and the invariant function
        \[ N_i := \prod_{l=1}^{n_i} t_i(x_l) = \varepsilon^{(i)}_{n_i} \]
        is a regular function near $z_\mdls$ vanishing exactly on $\Delta'_{P_i}$
        (as $t_i$ vanishes only at $P_i$ on $U_i$). In particular $w'_i$ divides
        $\varepsilon^{(i)}_{n_i}$ in $\Oo_{C^{(\deg\mdls)}, z_\mdls}$. Orders with respect to the
        maximal ideal are computed in the completion and are additive (the associated graded ring is a
        polynomial ring, hence a domain); since $\varepsilon^{(i)}_{n_i}$ is one of the coordinates of
        the power series ring, its order is $1$, and therefore
        $\operatorname{ord}_{z_\mdls}(w'_i) = 1$ as well. The last statement of the claim follows from
        \Cref{claim:exceptional_valuation_is_order}.
    \end{proof}

    \begin{claim}\label{claim:coefficients_sum_to_zero}
        $\sum_{i=1}^{s} c_i \equiv 0 \pmod{e}$.
    \end{claim}
    \begin{proof}
        Here we use the Abel--Jacobi morphism \emph{with trivial modulus}
        \[ \varphi_d : C^{(d)} \longrightarrow \PicC[d]. \]
        Since $d \geq 2g - 1 + \deg\mdls \geq 2g$ and $k$ is algebraically closed (so $C$ has a rational
        point), $C^{(d)}$ is a projective space bundle over $\PicC[d]$ and the fiber over a point
        $[\cL]$ is the complete linear system $|\cL| = \bP\left(H^0(C, \cL)\right)$
        (\Cref{section:AbelJacobi}).
        Let $\xi$ be the generic point of $\PicC[d]$, with residue field $\kappa(\xi)$, and let
        \[ P_\xi := \varphi_d^{-1}(\xi) \cong \bP^{\,d-g}_{\kappa(\xi)} \]
        be the generic fiber; it is a localization of $C^{(d)}$, so for every point
        $w \in P_\xi$ we have $\Oo_{P_\xi, w} = \Oo_{C^{(d)}, w}$. In particular the codimension-one
        points of $P_\xi$ are exactly the generic points of the prime divisors of $C^{(d)}$ that
        dominate $\PicC[d]$, and for such a divisor $W$ with generic-fiber components
        $w \in W \cap P_\xi$,
        \begin{equation}\label{equation:restriction_to_generic_fiber}
            \operatorname{div}_{P_\xi}\left( F|_{P_\xi} \right)
            = \sum_{W} v_W(F) \cdot \left[ W \cap P_\xi \right],
        \end{equation}
        where the sum is over the prime divisors of $C^{(d)}$ dominating $\PicC[d]$ and
        $[W \cap P_\xi]$ denotes the corresponding reduced cycle of codimension one on $P_\xi$.
        (The generic point of $C^{(d)}$ lies in $P_\xi$, so $F$ restricts to a nonzero rational function
        on $P_\xi$.)

        Each $\Delta_{P_i}$ dominates $\PicC[d]$: for every $\cL$ of degree $d$ one has
        $h^0(\cL(-P_i)) = d - g \geq 1$, so every fiber $|\cL|$ contains divisors through $P_i$.
        Moreover, inside the generic fiber,
        \[ \Delta_{P_i} \cap P_\xi = \left\{ D \in |\cL_\xi| \; : \; D \geq P_i \right\}
           = \bP\left( H^0(C_{\kappa(\xi)}, \cL_\xi(-P_i)) \right)
           \subset \bP\left( H^0(C_{\kappa(\xi)}, \cL_\xi) \right), \]
        a linear subspace of codimension one (as $\deg \cL_\xi(-P_i) = d - 1 \geq 2g - 1$, so
        $h^0(\cL_\xi(-P_i)) = d - g = h^0(\cL_\xi) - 1$). In particular $\Delta_{P_i} \cap P_\xi$ is
        irreducible and, as a reduced linear subvariety of projective space, it has degree $1$.

        Now take degrees in \eqref{equation:restriction_to_generic_fiber}. A principal divisor on
        $\bP^{\,d-g}_{\kappa(\xi)}$ has degree $0$ (its homogeneous numerator and denominator have equal
        degrees). By \Cref{claim:divisor_congruence}, every dominating prime divisor $W$ other than the
        $\Delta_{P_i}$ contributes a multiple of $e$, while $\Delta_{P_i}$ contributes
        $c_i \cdot \deg[\Delta_{P_i} \cap P_\xi] = c_i$. Hence
        \[ 0 = \deg \operatorname{div}_{P_\xi}\left( F|_{P_\xi} \right)
             \equiv \sum_{i=1}^{s} c_i \pmod{e}. \qedhere \]
    \end{proof}

    We can now conclude. Since $A$ is a unique factorization domain, we may write
    \[ F = u \cdot \prod_{j} w_j^{\,a_j}, \qquad u \in A^\times, \]
    where $w_j$ runs over generators of the height-one primes of $A$, i.e.\ over local equations at
    $\zeta$ of the prime divisors $W_j$ of $C^{(d)}$ passing through $\zeta$, and $a_j = v_{W_j}(F)$.
    Applying $v_{\eta_0}$ and using \Cref{claim:exceptional_valuation_is_order}:
    \[ v_{\eta_0}(F) = \sum_j v_{W_j}(F) \cdot \operatorname{ord}_\zeta(w_j). \]
    By \Cref{claim:divisor_congruence}, all terms with $W_j \notin \{\Delta_{P_1}, \dots, \Delta_{P_s}\}$
    are divisible by $e$; by \Cref{claim:boundary_divisors_have_multiplicity_one},
    $\operatorname{ord}_\zeta(w_i) = 1$ for the boundary divisors. Therefore, using
    \Cref{claim:coefficients_sum_to_zero},
    \[ v_{\eta_0}(F) \equiv \sum_{i=1}^{s} c_i \equiv 0 \pmod{e}, \]
    which is \eqref{equation:kummer_valuation_at_boundary}. As explained above, this shows
    $\chi(\rho'(I)) = \{1\}$ for every character $\chi$ of $G'$, hence $\rho'(I) = \{1\}$; combined
    with $\rho_p(I) = \{1\}$ this gives $\rho(I) = \{1\}$, i.e.\ $\cQ$ is unramified at $(H, \eta_0)$.
\end{proof}

\subsubsection*{Conclusion of the proof}

\begin{proof}[Proof of \Cref{lemma:ExtensionOfTamelyRamifiedSheaf}]
    Let $G = \Lambda^\times$ and let $\cP^{[d]}$ be the $G$-torsor on $U^{(d)}_{\text{\'et}}$
    corresponding to $\cF^{[d]}$ under \Cref{prop:torsor_module_equivalence}, with character
    $\rho : \pi_1^{et}(U^{(d)}, \bar{u}) \to G$.
    By \Cref{lemma:boundary_of_compactification}(3), the only codimension-one point of
    $\Cmod{d} \setminus U^{(d)}$ is $\eta_0$, so the hypothesis that $\cF^{[d]}$ is tamely ramified
    along $H$ means precisely that $\cP^{[d]}$ is tamely ramified over $\Cmod{d}$ at $(H, \eta_0)$
    (\Cref{definition:tamely_ramified_in_codim_1}).
    By \Cref{prop:tame_at_boundary_implies_unramified}, $\cP^{[d]}$ is then unramified at
    $(H, \eta_0)$. By \Cref{prop:extension_via_purity}, the character $\rho$ factors uniquely as
    $\rho = \tilde{\rho} \circ \iota_*$ through
    $\pi_1^{et}(\Cmod{d}, \bar{u})$. Let $\tilde{\cF}^{[d]}$ be the locally constant sheaf of free
    rank-one $\Lambda$-modules on $(\Cmod{d})_{\text{\'et}}$ corresponding to $\tilde{\rho}$
    (\Cref{cor:abelian_equivilance_fundamental_torsors} together with
    \Cref{prop:torsor_module_equivalence}). Since $\tilde{\rho} \circ \iota_* = \rho$, its restriction
    to $U^{(d)}$ is isomorphic to $\cF^{[d]}$.
\end{proof}

\begin{rem*}
    The proof also gives uniqueness: since $\iota_* : \pi_1^{et}(U^{(d)}) \to \pi_1^{et}(\Cmod{d})$ is
    surjective (\Cref{prop:extension_via_purity}), the extension $\tilde{\cF}^{[d]}$ is unique up to
    isomorphism, and more generally the restriction functor from locally constant sheaves on $\Cmod{d}$
    to locally constant sheaves on $U^{(d)}$ is fully faithful. This is the source of the uniqueness
    assertions in \Cref{thm:GCFT_torsors} and \Cref{thm:GCFT_reduced}.
\end{rem*}

\begin{rem*}
    Statements of this kind appear in both of the works our approach follows, with different proofs.
    Guignard (\cite{Guignard2018}, Lemmas 5.7 and 5.9) does not extend $\cF^{[d]}$ over a
    compactification: he shows that $\cF^{[d]}$ is constant on the (affine space) fibers of $\Phi_d$,
    using the triviality of the tame fundamental group of $\A^1$, and then descends it along $\Phi_d$
    directly by a duality argument. Takeuchi (\cite{takeuchi2019blow}, proof of Theorem 1.2 in \S 4)
    obtains the extension of the corresponding character over $\Cmod{d}$ from his refined
    (Kato--Matsuda) Swan conductor computation (\cite{takeuchi2019blow}, Theorem 2.7 and Corollary 2.8),
    which for $p$-power order characters shows unramifiedness at $\eta_0$ directly.
    The proof given here is instead elementary: beyond the purity of the branch locus, it only uses
    Kummer theory and the geometry of the linear systems $|\cL|$ inside $C^{(d)}$.
\end{rem*}
